pdfoutput=1
\pdfoutput=1
% !TEX program = xelatex
% ===========================================================================
%  SCX Business Gauge: Three-Layer Restructuring of Global AI Infrastructure
%  Commercial Landscape and Strategic Analysis
%  Date: 2026-07-01
% ===========================================================================
\documentclass[12pt,a4paper]{article}
% ---- Mathematics ----
\usepackage{amsmath,amssymb,amsthm,mathtools}
\usepackage{bm}
\usepackage{bbm}
\usepackage{dsfont}
% ---- Theorem Environments ----
\newtheorem{theorem}Theorem[section]
\newtheorem{proposition}[theorem]{Proposition}
\newtheorem{corollary}[theorem]Corollary
\newtheorem{lemma}[theorem]Lemma
\newtheorem{definition}[theorem]Definition
\theoremstyle{remark}
\newtheorem{remark}[theorem]Remark
\newtheorem{example}[theorem]{Example}
% ---- Layout ----
\usepackage[margin=2.5cm]{geometry}
\usepackage{booktabs}
\usepackage{graphicx}
\usepackage{hyperref}
\usepackage{enumitem}
\usepackage[capitalise,noabbrev]{cleveref}
\usepackage{xcolor}
\usepackage{tikz}
\usepackage{float}
\usepackage{longtable}
\usepackage{makecell}
\usepackage{multirow}
\usepackage{setspace}
\onehalfspacing
% ---- Hyperref ----
\hypersetup{
    colorlinks=true,
    citecolor=blue,
    urlcolor=blue,
    pdfauthor={SCX},
    pdfsubject={SCX Business Gauge},
}
% ---- Custom Commands ----
\newcommand{\SCX}{\textsf{SCX}}
\newcommand{\Yajie}{\textsf{Yajie}}
\newcommand{\Spring}{\textsf{Spring}}
\newcommand{\Situs}{\textsf{Situs}}
\newcommand{\Cercis}{\textsf{Cercis}}
\newcommand{\R}{\mathbb{R}}
\newcommand{\E}{\mathbb{E}}
\newcommand{\Pbb}{\mathbb{P}}
\newcommand{\Var}{\mathrm{Var}}
% Business analysis boxes
\newtheorem{businessnote}{Business Note}[section]
\renewcommand{\thebusinessnote}{\textcolor{blue}{\textbf{[Note \arabic{businessnote}]}}}
\newenvironment{businessbox}{%
  \begin{center}\color{blue}\begin{minipage}{0.92\textwidth}\begin{businessnote}}{%
  \end{businessnote}\end{minipage}\end{center}}
% Honest critique
\newtheorem{honesthit}{Honest Strike}[section]
\renewcommand{\thehonesthit}{\textcolor{red}{\textbf{[Strike \arabic{honesthit}]}}}
\newenvironment{hitbox}{%
  \begin{center}\color{red}\begin{minipage}{0.92\textwidth}\begin{honesthit}}{%
  \end{honesthit}\end{minipage}\end{center}}
% ---- Title ----
\title{\textbf{SCX Business Gauge: Three-Layer Restructuring of\\Global AI Infrastructure}\\[8pt]
       \large Protocol, Data/Compute, and Application Layers:\\Strategic Decoupling and Value Redistribution}\\[6pt]
       \normalsize --- With Company Valuations Under the Yajie Protocol}
\author{SCX}
\date{2026-7-1 \\ July 1, 2026}
\begin{document}
\maketitle
% ===================================================================
\begin{abstract}
\noindent
SCX is not a competitor in the AI industry. It is the infrastructure that reorganizes the AI industry. This paper presents a three-layer model of the post-SCX commercial landscape: \textbf{Layer 1 (Protocol)} --- SCX as the universal standard for data quality, enjoying a permanent monopoly through theorem-level non-circumventability; \textbf{Layer 2 (Data/Hardware)} --- DeepSeek, Alibaba, Huawei, NVIDIA providing the computational substrate, differentiable but substitutable; \textbf{Layer 3 (Application)} --- user-facing products on standardized infrastructure, characterized by thin margins and high churn.

We argue that SCX occupies the structural position of TCP/IP in the data quality stack: a protocol that nobody competes with because everyone builds on it. The protocol layer captures value through \textit{non-circumventability}---any audit requires citing SCX theorems, any \(M_t\) requires symbiotic binding, any Cercis score traces back to Yajie. We analyze winners and losers in the post-SCX equilibrium, present the ``Bible-and-Pope'' revenue model, and provide DCF-calibrated Yajie Valuation Multipliers (\(\mu_{\Yajie}\)) for major AI companies.

\medskip
\noindent\textbf{Keywords:} SCX, data quality infrastructure, three-layer model, protocol monopoly, non-circumventability, AI industry structure, business model, Bible-and-Pope, company valuation
\end{abstract}
\newpage
\tableofcontents
\newpage
% ===================================================================
\section{Introduction: The Three-Layer Revolution in AI}
\subsection{The AI Industry Before SCX}
Before SCX, the AI industry organized around a single axis: \textit{who has the best model}. Every major lab---OpenAI, Anthropic, Google DeepMind, DeepSeek, Meta AI---built models, curated datasets, developed evaluation frameworks, and shipped products within the same organizational boundary. This vertical integration obscured the fundamental tension: \textbf{model quality} and \textbf{data quality} are distinct problems requiring distinct organizational solutions.

\begin{businessbox}
The pre-SCX AI industry suffered from a classic \textbf{public goods underprovision problem}. Data quality assessment is inherently a public good: once a dataset is audited, the certification benefits all downstream consumers regardless of who paid for the audit. SCX solves this by transforming data quality from a private cost into a shared infrastructure asset.
\end{businessbox}

\subsection{The Three-Layer Architecture}
SCX introduces a structural decoupling reorganizing AI into three distinct layers:
\begin{enumerate}[label=\textbf{\arabic*}, leftmargin=*]
    \item \textbf{Protocol Layer (SCX)} --- Defines what ``correct'' means. Permanent monopoly via theorem non-circumventability.
    \item \textbf{Data/Hardware Layer} --- Provides compute and training data. DeepSeek, Alibaba, Huawei, NVIDIA.
    \item \textbf{Application Layer} --- User-facing products. Low margin, high churn.
\end{enumerate}

\subsection{The Core Thesis}
SCX does not compete in AI. It reorganizes AI. When a company says ``we compete with OpenAI,'' that is a statement about Layer 3. When they say ``we use SCX,'' that is about accepting a protocol standard---analogous to saying ``we use TCP/IP.''

\begin{hitbox}
The most dangerous strategic error is treating SCX as a competitor rather than as infrastructure. Companies that view SCX as ``another audit tool'' make the same mistake companies made with TCP/IP in the 1980s---attempting proprietary alternatives to a protocol that becomes universal \textit{because} it is a protocol.
\end{hitbox}
% ===================================================================
\section{Protocol Layer Economics}
\subsection{Theorem Non-Circumventability}
The SCX protocol layer derives monopoly power from \textbf{mathematical theorems}:
\begin{enumerate}[label=(\arabic*), leftmargin=*]
    \item \textbf{Theorem 1 (Noise Detection):} Multi-expert consistency achieves exponential F1 convergence. Any audit not using this \textit{cannot} achieve this guarantee.
    \item \textbf{Theorem 3 (Honest Person):} Distinguishing label noise from sample difficulty is \textit{mathematically impossible} from observational data alone.
    \item \textbf{Theorem 4 (Minimax Optimality):} SCX's adaptive threshold achieves the theoretical lower bound. No algorithm can achieve a smaller error constant---a permanent moat.
\end{enumerate}

\begin{definition}[Non-Circumventability]
A protocol is \textbf{non-circumventable} if any entity performing the function the protocol addresses \textit{must}, by mathematical necessity, operate within the protocol's framework.
\end{definition}

\subsection{Network Effects Without Users}
SCX exhibits \textbf{network effects without users}. Its value is tied to \textbf{cumulative audit volume}---the number of datasets processed through the system---not user count. The protocol's accuracy \(\theta(t)\) improves as:
\begin{equation}
\theta(t) = \theta_{\max} - (\theta_{\max} - \theta_0) e^{-\gamma \cdot |\mathcal{E}_t|}
\end{equation}
where \(|\mathcal{E}_t|\) grows with \textit{time and usage volume}, not \textit{user count}.

\subsection{The M=1 UNDECLARED Problem}
Every AI lab faces a binary choice: adopt SCX (\(M \geq 8\) experts, verifiable quality) or remain \textbf{M=1 UNDECLARED} (single model, unverifiable claims). By Theorem 3, M=1 UNDECLARED claims are mathematically unverifiable.

\begin{hitbox}
\textbf{The ``trust me'' era of AI is over.} Companies claiming ``our data is carefully curated'' without SCX audit will discover these claims are worth exactly nothing in a market where competitors can \textit{prove} their data quality.
\end{hitbox}
% ===================================================================
\section{Data and Hardware Layer}
\subsection{The Strategic Pivot}
Post-SCX, a company's AI reputation fragments. DeepSeek's sustainable differentiation reduces to \textbf{GPUs and data volume}. Model architecture advantages diffuse quickly; training techniques are copied; evaluation benchmarks are gamed. The only non-replicable assets are the physical GPU fleet and proprietary datasets.

\subsection{The Bible-and-Pope Revenue Model}
\begin{itemize}
    \item \textbf{The Bible (Spring + Yajie Core):} Open-source, freely available protocol specification, theorems, and core algorithms.
    \item \textbf{The Pope (Yajie API Calibration):} Paid, production-grade audit services with guaranteed precision, CEC access, and continuous Cercis score updates.
\end{itemize}

Revenue streams: Pay-per-Audit, Continuous Monitoring, CEC Access (data license), Managed Infrastructure, SCX Trust Mark certification. The open-source Bible drives adoption; the calibrated Pope captures value through precision.

\subsection{Why Competitors Cannot Replicate}
Three insurmountable barriers: (1) Theorem originality---citations anchor SCX as the originator; (2) CEC time deficit---competitors start with \(\mathcal{E}_0 = 0\); (3) Open-source paradox---any competitor must explain why their system is more trustworthy than an open-source, theorem-backed protocol.
% ===================================================================
\section{Company Valuations Under the Yajie Protocol}
\subsection{The Yajie Valuation Multiplier \(\mu_{\Yajie}\)}
\begin{definition}[Yajie Valuation Multiplier]
\begin{equation}
\mu_{\Yajie} = \frac{V_{\text{post}}}{V_{\text{pre}}} 
= \alpha \cdot \mathbbm{1}_{\text{Layer 1}}
+ \beta \cdot \frac{M_{\text{declared}}}{M_{\min}}
+ \gamma \cdot \frac{|\mathcal{E}_t^{\text{company}}|}{|\mathcal{E}_t^{\SCX}|}
- \delta \cdot \mathbbm{1}_{\text{M=1 UNDECLARED}}
\end{equation}
where \(M_{\min}=8\), and \(\alpha, \beta, \gamma, \delta > 0\) are market-calibrated weights.
\end{definition}

\subsection{DCF-Calibrated \(\mu_{\Yajie}\) by Company}
\begin{longtable}{@{}p{2.2cm} p{1.8cm} p{1.6cm} p{1.8cm} p{2cm} p{2cm}@{}}
\caption{DCF-Calibrated \(\mu_{\Yajie}\) Ranges}
\toprule
\textbf{Company} & \textbf{\(V_{\text{pre}}\) (\$B)} & \textbf{Layer} & \textbf{UNDEC?} & \textbf{\(\mu_{\Yajie}\)} & \textbf{\(V_{\text{post}}\) (\$B)} \\
\midrule
SCX & N/A & L1 & No & 3.0--8.0\(\times\) & \(\sim\)50B+ \\
NVIDIA & \$2,800 & L2 & No & 0.65--0.85 & \$1,820--\$2,380 \\
Huawei (Ascend) & \$150--300 & L2 & No & 1.5--2.5 & \$225--\$750 \\
DeepSeek & \$20--50 & L2 & No & 2.0--4.0 & \$40--\$200 \\
OpenAI & \$150 & L2/L3 & \textbf{Yes} & 0.3--0.6 & \$45--\$90 \\
Anthropic & \$60 & L2/L3 & \textbf{Yes} & 0.3--0.5 & \$18--\$30 \\
Google DeepMind & \$2,100 & L2/L3 & Partial & 1.1--1.4 & \$2,310--\$2,940 \\
Meta AI & \$1,300 & L2/L3 & Partial & 0.8--1.0 & \$1,040--\$1,300 \\
xAI (Grok) & \$40 & L3 & \textbf{Yes} & 0.4--0.6 & \$16--\$24 \\
\bottomrule
\end{longtable}

\subsection{NVIDIA: Short-Term Bull, Long-Term Bear}
\textbf{Short-term (2026--2028):} Yajie's \(M \geq 8\) requirement multiplies GPU demand by 6--8\(\times\) for the same model generation. \textbf{Long-term (2029+):} Spring noise detection reduces training data requirements; Yajie commoditizes model quality, making GPU brand less relevant; CUDA's software moat erodes under framework abstraction.

\subsection{Structural Winners and Losers}
\textbf{Winners:} Independent researchers (equal audit footing), data providers (premium pricing for audited data), hardware manufacturers (\(M\times\) GPU multiplier), cloud providers (managed audit infrastructure), open-source communities (Spring is open).

\textbf{Losers:} Closed-source model companies (cannot prove quality), proprietary audit services (Yajie is free and better), M=1 UNDECLARED labs (trust premium evaporates), unaudited data brokers (information asymmetry collapses).
% ===================================================================
\section{Conclusion}
SCX is the TCP/IP of the AI data quality stack. It does not compete---it reorganizes. The three-layer architecture (Protocol / Data-Hardware / Application) creates a structural bifurcation: audited systems that can prove their quality, and unaudited systems whose claims are mathematically unverifiable. The Bible-and-Pope model monetizes precision, not access. The Yajie Valuation Multiplier provides a quantitative framework for repricing AI assets in the post-SCX equilibrium. The transition is inevitable---the only question is timing.
% ===================================================================
\begin{thebibliography}{99}
\bibitem{scx2026theorems} SCX. \emph{SCX Core Theorems}. Technical Report, 2026.
\bibitem{scx_galois} SCX. \emph{Galois-SCX}. Preprint, 2026.
\bibitem{scx_compactness} SCX. \emph{Compactness-SCX}. Preprint, 2026.
\end{thebibliography}
\end{document}
